\subsection{AI Explainability and Transparency}
\label{sec:6.3.4}
An AI system entrusted with decisions needs to be understandable, not just correct: explainability and transparency give developers, users, regulators, and the people affected by a decision an account to check the decision against. Section~\ref{sec:6.3.3} states the limit that governs everything below: an account of a decision can be produced after it, so what these techniques deliver is a decision that can be contested rather than a window onto what produced it. As AI algorithms grow more complex, that inspection gets harder exactly as the decisions being made get more consequential — which is what makes explainability and transparency safety infrastructure.

Explainability is the ability to clarify the rationale behind a specific decision or output. If a medical-diagnostic AI system recommends a particular treatment, understanding how it reached that conclusion lets a doctor evaluate the recommendation on its merits and catch a bias or inaccuracy in the system's data or algorithm before it reaches the patient. Transparency is the wider practice of sharing an AI system's inner workings — its algorithms, training data, objectives — so a stakeholder can confirm the system adheres to the standards it claims to, rather than taking the claim on faith: a recidivism-risk system used in criminal justice that discloses how it weighs its inputs makes it possible to catch and correct bias against a demographic group, in a way a sealed system does not.

Achieving either at a level that satisfies everyone involved runs into limits. A deep neural network is a “black box” in the literal sense that its internal representations resist human interpretation, and full transparency carries its own risk: exposing an AI system's training data or objectives can compromise privacy or intellectual property, or hand a bad actor a map of exactly how to exploit it. The European Union's General Data Protection Regulation addresses part of this by requiring an explanation of the logic behind automated decisions that significantly affect an individual \autocite{europeanunion2016general}, and interdisciplinary work between computer scientists, ethicists, and policymakers continues on the rest.

\runin{Challenges to Explainability and Transparency}
\begin{itemize}
  \item Complexity and nonlinearity: a deep learning architecture's many layers, nonlinear relationships, and interactions among variables resist interpretation by construction, not as an oversight that better tooling can simply remove.
  \item Accuracy versus explainability: a model tuned to also be transparent sometimes gives up some of the performance a less interpretable model would have delivered, and there is no general method for knowing in advance how much.
  \item Identifying which features matter: even a comparatively simple, transparent model can leave it unclear which inputs drove a given decision, especially with high-dimensional data.
  \item Diverging standards: designers, users, and regulators disagree about what counts as an adequate explanation, and no standard has emerged that satisfies all three at once.
  \item Technique mismatch: an explainability method built for one kind of model does not necessarily transfer to another, so novel architectures can outrun the tools built to explain them.
  \item Legal and ethical exposure: explaining an AI system's reasoning can require revealing information — proprietary methods, training data that touches personal information — that the organization has other obligations to protect.
\end{itemize}
\runin{Explainability Techniques}
LIME and SHAP are the two tools practitioners reach for, and section~\ref{sec:6.1.3} sets out what each does and where both stop: neither can tell you the model is aimed at the wrong thing, because both take the training target as given.

Decision trees and rule-based systems trade some of the predictive power of a black-box model for interpretability that comes built in: a decision can be read directly off a sequence of human-readable rules or tree branches, which matters in domains like finance, medicine, or law enforcement, where the reasoning has to survive scrutiny on its own terms and not just the outcome.

Counterfactual explanation asks what would have had to be different: in a credit-approval system, identifying the smallest change to an applicant's financial profile that would have flipped the decision tells the applicant what mattered and gives the institution a check on whether that factor should have mattered as much as it did.

Visualization and interactive tooling make an AI system's behavior something a person can explore. Dimensionality-reduction techniques like t-SNE and UMAP turn high-dimensional data into something visually inspectable, and interactive tools that let a user vary an input and watch how a model's output responds make the model's behavior tangible in a way a static explanation cannot.
